\documentclass[11pt]{article}

\usepackage[T1]{fontenc}
\usepackage{lmodern}
\usepackage[strings]{underscore}
\usepackage[letterpaper,margin=1in]{geometry}
\usepackage{microtype}
\usepackage{parskip}
\usepackage{array}
\usepackage{booktabs}
\usepackage{enumitem}
\usepackage[hidelinks]{hyperref}

\newcolumntype{L}[1]{>{\raggedright\arraybackslash}p{#1}}

\setlist[itemize]{topsep=4pt, itemsep=4pt, parsep=3pt}
\setlist[enumerate]{topsep=4pt, itemsep=6pt, parsep=3pt}

\title{Proving Liveness of Load-Store Queue}
\author{}
\date{}

\begin{document}
\maketitle

Follow the \emph{relational ranking} approach of McMillan, \emph{Toward
Liveness Proofs at Scale}, CAV 2024 (Rules~6 and~8). Fix \texttt{e} to be
an arbitrary entry of the load queue, at slot \texttt{K}.

There needs to be some sort of statement that infinitely often the LSQ accepts
requests. 



\section{In queue entails will be sent}
Rule 6, instantiations:
\begin{itemize}
\item \texttt{p}: \texttt{e} is in the queue --- \texttt{load_pending(K)}
\item \texttt{q}: \texttt{e} is answered in the queue --- \texttt{load_rdy(K)}
\item \texttt{r}: The memory system responds to an entry before \texttt{e}
\item Liveness invariant \texttt{phi}: \texttt{true}. It isn't needed.
\item Ranking relation \texttt{delta}: the things ahead of \texttt{e}, in
  \emph{both} queues, that are not ready --- the \emph{set}, not its
  cardinality.
\end{itemize}

This yields the following 4 invariants to prove:
\begin{itemize}
\item 
\end{itemize}

\subsection{The ranking relation}

For a pending load slot \texttt{K}, with \texttt{off(X) = X - load_head},
\texttt{delta(K)} has one arm per queue:

{\small
\begin{verbatim}
loads_not_ready_before(K, L) =
    load_pending(L) & ~load_rdy(L) & off(L) < off(K)

stores_not_ready_before(K, J) =
    store_pending(J) & mark_offset(J) <= off(K)
\end{verbatim}
}

\texttt{delta(K)} is the pair, ordered componentwise: it is
\emph{conserved} when neither arm gains an element and \emph{reduced}
when one arm loses one and neither gains. Two arms rather than one
because the two queues are indexed by different sorts, \texttt{load_idx}
and \texttt{store_idx}; the paper's Definition~1 admits exactly this ---
a ranking is an indexed set of relations, ``each predicate may be over a
different sort''. This is the same split the two summands of a numeric
measure would have had, minus the addition.

\texttt{mark_offset(J) <= off(K)} is ``store \texttt{J} is older than load
\texttt{K}'' in pure offset arithmetic (\texttt{cross_s}, \texttt{cross_l})
--- the same idiom \texttt{store_elig} and
\texttt{l_gap} already use, so no new tag reasoning enters. A pending
store is never \emph{ready}; it leaves only by departing, so every older
pending store counts, which is what makes the second arm the store
queue's share of the ranking.

Both arms are indexed off a register (\texttt{load_pending}/\texttt{store_pending}
compare against \texttt{load_head}/\texttt{store_head}), so no unrolling
over the four offsets is needed and no \texttt{while} appears: as
relations they are stated once, universally, and the solver instantiates
them at the slots it needs. This is the second thing the relational form
buys --- the numeric version had to be written out four times per arm to
be summable.

\subsection{Reused relations}

The ranking needs three things and two of them already exist.

\begin{itemize}
\item \texttt{stores_not_ready_before(K, store_head)} \textbf{is}
  \texttt{\textasciitilde{}load_elig(K)} for a
  pending \texttt{K}: \texttt{load_elig}'s body says either the store
  queue is empty --- so \texttt{store_pending(store_head)} is false ---
  or \texttt{store_head}'s mark sits strictly above \texttt{off(K)}.
  Worth stating once, since statements~1 and~3 both lean on it:

  {\small
  \begin{verbatim}
invariant [elig_zero] load_pending(K) ->
    (load_elig(K) <-> ~stores_not_ready_before(K, store_head))
\end{verbatim}
  }

  It is also the bridge that lets the store arm inherit the existing
  tag-order facts (\texttt{cross_s}, \texttt{cross_l}) rather than
  needing new ones. Note what is \emph{not} needed: the numeric
  version had to turn ``the counted set is non-empty'' into
  ``\texttt{store_head} is in it'', which is \texttt{m_mono}. A
  relation is reduced by exhibiting a
  witness, and statement~3 names \texttt{st_head} as the witness in its
  own hypothesis, so \texttt{m_mono} no longer has to carry that step.
\item \texttt{store_pending} and \texttt{load_pending} are the window
  guards in the two arms; \texttt{load_pending_pre} is just the
  latter's body restated over the pointer wires.
\item \texttt{store_elig} is \emph{not} reusable for the
  load arm, though it reads close to it. Its body lets a load ahead be
  sent-but-unanswered provided the addresses differ; the ranking has to
  include that load, because it is not ready and the cache still owes
  it. So \texttt{loads_not_ready_before} is the one genuinely new
  relation here.
\end{itemize}

Why this ranking and not ``the entries before \texttt{e}'':

{\small
\begin{tabular}{@{} L{0.34\linewidth} L{0.60\linewidth} @{}}
\toprule
event & effect on \texttt{delta(K)} \\
\midrule
response captured at an offset \texttt{< off(K)} & \textbf{loses \texttt{cache_to_lsq_resp_idx}} \\
response captured at \texttt{K} & \texttt{q} reached \\
response captured after \texttt{K} & conserved \\
store departs (\texttt{st_go}) and was older than \texttt{K} & \textbf{loses \texttt{st_head}} \\
store arrives & conserved --- its mark lands at \texttt{off(ld_tail)}, younger than \texttt{K} \\
load arrives & conserved --- it lands at the tail \\
head load retires (\texttt{ld_ans}) & conserved --- the retired head was ready, so it was never in \texttt{delta} \\
port grant to a load & conserved \\
\bottomrule
\end{tabular}
\par}

The plain ``entries before \texttt{e}'' relation loses an element only on
a retire, which is not what the cache's justice condition delivers; this
one loses one on exactly the two events the cache is fair about.

Note the retire row. Both arms are stated in \emph{offsets} from
\texttt{load_head}, and a retire shifts \texttt{load_head} and
\texttt{mark_offset} together, so every membership is
preserved across the shift. As a set of \emph{slots} the extension does
not move at all, which is what ``conserved'' has to mean.

\section{The four statements}

\begin{enumerate}
\item \textbf{Enabledness at the bottom of the ranking.} In the rule
  this premise is \texttt{phi -> En(tau)}, a \emph{single-state} formula
  --- note the absence of any \texttt{_pre} --- so it is a plain
  \texttt{invariant} in \texttt{implementation} next to \texttt{outst}
  and \texttt{pres_l}, evaluated in one state rather than paired across
  an edge, and it needs no sample and no new wire:

  {\small
  \begin{verbatim}
load_pending(K) & load_elig(K)
    & (forall L. ~loads_not_ready_before(K, L))
    & ~load_rdy(K) & ~load_sent(K) & ~req_en_r
    -> <the arbiter presents K on the port>
\end{verbatim}
  }

  Written with \texttt{load_elig} rather than an empty store arm: by
  \texttt{elig_zero} those say the same thing, and \texttt{load_elig} is
  the vocabulary the arbiter's own guard, \texttt{pres_l}, and
  \texttt{outst} are already stated in. What remains, the empty load
  arm, is the whole content of the statement: every slot ahead of
  \texttt{K} is ready, so none of them satisfies the arbiter's guard,
  so \texttt{K} wins the priority chain --- which resolves
  lowest-offset-last, hence lowest-offset-wins.

  The bottom of the ranking is needed here. Above it the helpful
  transition is not \texttt{K}'s at all but whatever removes an element
  ahead of it, and \texttt{K}'s own request is genuinely blocked:
  \texttt{load_elig(K)} fails while an older store is pending, and a
  lower-offset unready load outranks \texttt{K}.
  Justice on \texttt{creq_ready} then takes \texttt{K} off the port.

\item \textbf{The ranking is conserved.} This is \texttt{conserves delta},
  i.e.\ \texttt{forall x. delta'(x) -> delta(x)}, one conjunct per arm:

  {\small
  \begin{verbatim}
load_pending_pre(K)
    -> load_rdy(K)
     | ( (loads_not_ready_before(K, L)  -> loads_not_ready_before_pre(K, L))
       & (stores_not_ready_before(K, J) -> stores_not_ready_before_pre(K, J)) )
\end{verbatim}
  }

  No comparison and no arithmetic on the ranking: the unprimed
  relations are read post-edge and the \texttt{_pre} ones are the
  pre-edge samples. The \texttt{load_rdy} disjunct absorbs the retire
  case, where \texttt{off(K)} wraps to 7. It is sound there: the
  retiring slot cannot be re-enqueued on the same edge, because
  \texttt{cpu_ready_r} is low whenever the queue is already at capacity
  (\texttt{q_l}, \texttt{q_rdy}), so \texttt{load_rdy(K)} still holds
  post-edge.

\item \textbf{The helpful step reduces it.} \texttt{reduces delta} is
  \texttt{exists x. delta(x) \& \textasciitilde{}delta'(x)}, and in both
  halves the witness is already available, so it can be named rather
  than quantified. Two halves, because \texttt{cache_to_lsq_resp_valid}
  on its own reduces nothing --- the response may be for a slot after
  \texttt{K}, or match no sent entry at all; nothing in \texttt{lsq_body}
  constrains \texttt{resp_idx}, the whole channel is a free input.

  {\small
  \begin{verbatim}
load_pending_pre(K) & ld_capture_pre
    & (resp_idx_pre - lh_pre) < (K - lh_pre)
    -> loads_not_ready_before_pre(K, resp_idx_pre)
     & ~loads_not_ready_before(K, resp_idx_pre)

load_pending_pre(K) & st_go_pre & stores_not_ready_before_pre(K, sh_pre)
    -> ~stores_not_ready_before(K, sh_pre)
\end{verbatim}
  }

  The load half's witness is \texttt{resp_idx_pre}, the response index
  sampled at the start of the edge: whichever of \texttt{ld_capture}'s
  four per-offset guards fired, that slot is exactly the one the
  sample recorded, so the existential witness is already in hand.
  \texttt{ld_capture} itself is new only as the disjunction of those
  four guards, each already the datapath's own capture condition,
  restated.

  The store half's hypothesis,
  \texttt{stores_not_ready_before_pre(K, sh_pre)}, is the pre-edge form
  of \texttt{\textasciitilde{}load_elig(K)} (\texttt{elig_zero} again)
  and is simultaneously the witness the conclusion discharges. It is
  carried as a sample because \texttt{load_elig} reads
  \texttt{mark_offset}, which by the time this invariant is checked a
  same-edge retire may already have shifted.

\item \textbf{Justice.} An assumption on the environment, with nothing
  to prove: the cache accepts the port (\texttt{cache_to_lsq_creq_ready})
  infinitely often, and answers every load it has taken. Documented as
  a comment where the freeness of the cache channel is already
  described. This is stronger than ``the cache responds infinitely
  often'': the capture guard tests \texttt{resp_idx} \emph{and} the
  echoed address, so bare \texttt{resp_valid} liveness would leave a
  sent load unanswered forever.
\end{enumerate}

\section{Soundness}

\texttt{delta} is a relation over \texttt{load_idx} and \texttt{store_idx},
which \texttt{basis} interprets as \texttt{bv[3]}.
Inclusion on subsets of a finite set is well-founded, so no order needs
axiomatising and no induction axiom needs instantiating --- which is the
property the paper is buying, and the reason it replaces the numeric
rank.

Finiteness, premises F1/F2 of Rule~5, is therefore free: the paper needs
that rule only for rankings over infinite sorts, and it discharges it by
showing each transition adds boundedly many elements. Here the sorts are
finite by construction, so \texttt{delta} is finite in every state
whatever it contains. Even taking the obligation at face value it is
immediate: one edge enqueues at most one entry, a load or a store, never
both --- the accept step is guarded by \texttt{lsq_to_cpu_ready} and
branches on \texttt{cpu_to_lsq_req_wen} --- so \texttt{delta} can gain
at most that one entry --- and the offset
guards exclude it, since a new entry lands at the tail, younger than the
pending \texttt{K}. That is the row-by-row content of the event table
above, and it is what makes statement~2 provable rather than merely
true.

The argument the four statements feed is Rule~6, with the two halves of
statement~3 standing for the two justice conditions of Rule~8 under the
priority scheduler the paper describes: \texttt{psi_store} is ``an older
store is still pending'', i.e.\ \texttt{stores_not_ready_before(K, st_head)},
and \texttt{psi_load} is its negation. Both are stable in the paper's
sense --- a pending store leaves only by departing (\texttt{st_go}),
and once no older store is pending none can appear, since a new store's
mark lands at \texttt{off(ld_tail)} --- and their disjunction is
\texttt{true}, so some justice condition is always scheduled.
Statement~1 discharges the empty ranking, statement~3 says a non-empty
one is left strictly smaller by the scheduled condition, statement~2
says nothing in between undoes that, and statement~4 is what forces the
reducing step to actually be taken rather than merely stay possible.

\section{Not covered}

This proves \texttt{load_pending(e) -> eventually load_rdy(e)}. Reaching
\emph{``the cpu sees the response''} additionally needs
\texttt{load_rdy(e) -> eventually ld_ans} at \texttt{e}, which is the
``entries before \texttt{e}'' relation with the retire as the helpful
step --- a second, much smaller ranking argument on top of this one.

\end{document}
